\documentclass[10pt,a4paper]{amsart}
\usepackage[margin=19mm]{geometry}
\usepackage{amsmath,amssymb,mathtools}
\usepackage[scr=rsfs,cal=euler]{mathalfa}
\usepackage{xfrac}
\usepackage[unicode,hidelinks,bookmarks=false]{hyperref}
\newcommand{\ie}{i.e.\ }
\newcommand{\cf}{cf.\ }
\newcommand{\resp}{respectively\ }
\newcommand{\Subsection}{\subsection}
\providecommand{\nobreakdash}{\nobreak\hbox{-}}
\setlength{\emergencystretch}{2em}
\multlinegap0pt
\usepackage{bm}
\usepackage{bbm}
\usepackage{dsfont}
\usepackage{xspace}
\usepackage{cancel}
\usepackage{mathtools}
\usepackage{amsthm}
\makeatletter
\@ifundefined{theorem}{\newtheorem{theorem}{Theorem}}{}
\@ifundefined{lemma}{\newtheorem{lemma}[theorem]{Lemma}}{}
\makeatother
\title[On a theorem of Nikol'skii]{On a theorem of Nikol'skii}
\author{Asuman G. Aksoy, Qidi Peng}
\date{Source: arXiv:2509.16344v1 (CC BY 4.0)}
\begin{document}
\maketitle

\noindent\emph{Source and licence.} On a theorem of Nikol'skii. Version arXiv:2509.16344v1 (CC BY 4.0), distributed under the Creative Commons Attribution 4.0 International licence, as recorded in the arXiv licence metadata for this exact version and shown on the abstract page https://arxiv.org/abs/2509.16344v1. Full rights notice: licenses/arxiv-2509.16344-CC-BY-4.0.txt.\par\smallskip

\noindent\textit{Editorial scope.} Counted unit: the Theorem whose source label is \texttt{Borodin} in the article (source0.tex lines 476--494, source label \texttt{Borodin}), delivered together with the article's own complete original proof of it (source0.tex lines 495--763, one \texttt{proof} environment of 269 lines). No mathematics is written, repaired or re-derived: statement and proof are the article's own text line for line. Required context retained uncounted: lemma1 (lines 116--120); lemma1' (lines 429--441); lemma2 (lines 446--471). Every definition, display and label of those lines is retained unchanged. Omitted from this package and not redistributed: 1--115, 121--428, 442--445, 472--475, 764--1220. Nothing inside a retained excerpt is omitted or altered: every retained body is delivered exactly as the article's lines are written, so no omission receipt of any kind accompanies this package. Citations: the retained excerpts carry no citation command at all, so no bibliography entry of the article is transcribed and no reference is redistributed. Reference adaptation: none was needed and none was made; every reference inside the delivered unit resolves inside it, so no unresolved reference or citation is produced. Printed slips kept literal: no printed word, symbol, hypothesis or number of the counted statement or of its proof was altered. Layout and class: the article's own class is \texttt{amsart} with packages including amsmath, amsthm, amscd, amsfonts, amssymb, graphicx, color, hyperref, float, lineno, epstopdf, bbm, mathrsfs, dsfont; the wrapper is the lane's pinned \texttt{amsart} editorial wrapper at 10pt on A4, which restates the theorem-like environments the retained text needs and adds the article's own macro definitions that the retained text needs (none), with extra wrapper packages (bm, bbm, dsfont, xspace, cancel, mathtools). The delivered numbering is the unit's own. Assets: no retained span contains a figure environment, a graphics-inclusion command, a PDF-page inclusion, a table or a TikZ picture; no image file is copied into this package, the package declares no asset dependency and no image file is redistributed with it. Licence: version arXiv:2509.16344v1 of this article is distributed under the Creative Commons Attribution 4.0 International licence, as recorded in the arXiv licence metadata for this exact version and shown on the abstract page https://arxiv.org/abs/2509.16344v1. A full rights notice is delivered at licenses/arxiv-2509.16344-CC-BY-4.0.txt. Characters: every retained excerpt is pure ASCII, so the delivered engine, XeLaTeX with its default Latin Modern text font, reports no missing character.
\begin{lemma}
\label{lemma1}
Let $(X,\|\cdot\|)$ be a normed linear space, $Y_1\subset Y_2\subset\ldots\subset Y_n\subset X$ be a finite system of strictly nested subspaces, $d_1>d_2>\ldots>d_n\ge0$ and $z \in X\backslash Y_n$. Then, there is an element $x\in X$ for which $\rho(x,Y_k)=d_k$ $(k=1,\ldots,n)$,  $\|x\|\le d_1+1$, and $x-\lambda z\in Y_n$ for some $\lambda>0$. If all $Y_k$ are existence subspaces (in particular, if all $Y_k$ are finite-dimensional or $X$ is reflexive),
then the above assertion holds for arbitrary numbers $d_1 \ge d_2 \ge \ldots\ge d_n \ge 0$.
\end{lemma}

\begin{lemma}
\label{lemma1'}
Let $(X,\|\cdot\|)$ be a normed linear space and let $Q$ be a subspace of $X$ with $\overline Q\subset X$. Pick two elements $x_1,x_2\in X\backslash \overline Q$ such that $x_2\notin \mbox{span} [\{x_1\}\cup Q]$. Then, there exists a non-zero linear functional $f:~X\longmapsto \mathbb R$ such that
\begin{equation}
\label{linearf1}
f(q)=0,~\mbox{for all $q\in Q$},~\|f\|=\frac{1}{\rho(x_1,Q)},~f(x_1)=1
\end{equation}
and
\begin{equation}
\label{linearf2}
f(x_2)=\lim_{a\to+\infty}\left(a-\frac{\rho(x_2-a x_1,Q)}{\rho(x_1,Q)}\right).
\end{equation}
\end{lemma}

 \begin{lemma}
 \label{lemma2}
 Let $Q_1,Q_2$ be two subspaces of an arbitrary normed linear space $(Q_3,\|\cdot\|)$, such that $\overline Q_k\subset Q_{k+1}$ for $k=1,2$. Also assume that for $k=1,2$ there is $y_k\in Q_{k+1}\backslash  Q_k$ for which $\rho(y_k,Q_k)=\|y_k\|$. Let $\{u_m\}_{m\ge1}$ and $\{v_m\}_{m\ge1}$ be two sequences of non-negative numbers, with $u_m\ge v_m$ for all $m\ge1$. Then there exists a sequence of elements  $\{q_m\}_{m\ge1}\subset Q_{3}$, given by
$$
q_{m}=v_m(z-f(z)w)+\mu_mf(z)w,~\mbox{for $m\ge1$},
$$
such that
\begin{enumerate}
\renewcommand{\labelenumi}{\alph{enumi})}
\item [(i)]$z$ is an element in $Q_3\backslash \overline Q_2$ satisfying $\rho(z,Q_1)=\rho(z,Q_2)=1$; $w$ is an element in $\overline Q_3\backslash\{0\}$ such that $\|z-w\|=1$.
    \item[(ii)] $f:~Q_3\rightarrow \mathbb R$ is a non-zero linear functional such that 
    $$
    f(z)=\lim_{a\to+\infty}\left(a-\frac{\rho(z-aw,Q_1)}{\rho(w,Q_1)}\right).
    $$
    \item[(iii)] For each integer $m\ge1$, the variable $\mu_m$ in the expression of $q_m$ is some value in $[v_m,u_m]$ such that
     \begin{equation}
 \label{q}
 \rho(q_m,Q_1)=u_m,~\rho(q_m,Q_2)=v_m,~\mbox{for all $m\ge1$}.
 \end{equation}
    \item[(iv)] The following inequalities hold:
 \begin{equation}
 \label{diffq1}
 \|q_m-q_n\|\le (\|z\|+2)\left(\max\{u_m,u_n\}-\min\{v_m,v_n\}\right),~\mbox{for all $m,n\ge1$}.
 \end{equation}
\end{enumerate}
 \end{lemma}

\begin{theorem} \label{Borodin}
Let $X$ be an arbitrary infinite-dimensional Banach space, $\{Y_n\}_{n\ge1}$ be a system of strictly nested subspaces with the properties
\begin{enumerate}
\item $\overline Y_n \subset Y_{n+1}$ for all $n\ge1$;
\item For each $n\ge1$, there exists $y_n\in Y_{n+1}\backslash Y_n$ such that
\begin{equation}
\label{Yn}
\rho(y_n,Y_n)=\|y_n\|.
\end{equation}
 \end{enumerate}
 Let  $\{d_n\}_{n\ge1}$ be a sequence of non-negative numbers  decreasing to $0$:
 \begin{equation}
\label{condition'}
 d_n\geq d_{n+1}\ge0, \,\,\mbox{for}~ n \ge 1;~\lim_{n\to\infty}d_n=0 \\
 % \sum_{k = n+1}^\infty d_k~\mbox{for all $n \ge n_0$ at which $d_n > 0$.
 \end{equation}
 and  $$ d_n > \sum_{k = n+1}^\infty d_k~\mbox{for all}\quad n \ge n_0$$
 Then there exists an element $x\in X$ such that $\rho(x,Y_n)=d_n$ for all $n\ge1$ 
\end{theorem}

\begin{proof} We start by assuming that $Y_1 \neq \{0\}$ and  $d_n>0$, for all $n\ge1$. The case $Y_1=\{0\}$ and the case when $d_n=0$ for some $n$ will be discussed separately in the end of this proof.

We define the sequence $\{\tau_j\}_{j\ge1}$ to be
$$
\tau_1=d_1;~\tau_{j}=\min\limits_{k=2,\ldots,j}\{d_{k-1}-d_{k}\},~\mbox{for $j\ge2$}.
$$
In view of the assumption on $\{d_n\}$, we know that the sequence $\{\tau_{j}\}_{j\ge1}$ is non-negative and decreasing to $0$. Since $Y_j\ne \{0\}$, then we can define $Y_0=\{0\}$ and for any integers $j$, $n$ with $1\le j\le n$, we can set
$$
u_n^{(j)}=1+\frac{\tau_n}{2^jd_j},~v_n^{(j)}=1.
$$
\begin{itemize}
\item For $j=1$, in view of Lemma \ref{lemma2}, we can find a sequence $\{q_{1,n}\}_{n\ge 1}\subset Y_{2}\backslash Y_1$ defined by
\begin{eqnarray*}
q_{1,n}&=&v_n^{(1)}(z_1-f_1(z_1)w_1)+\mu_n^{(1)}f_1(z_1)w_1\\
&=&(z_1-f_1(z_1)w_1)+\mu_n^{(1)}f_1(z_1)w_1,
\end{eqnarray*}
such that
\renewcommand{\labelenumi}{\alph{enumi})}
\begin{enumerate}
\item[(i)] $z_1$ is an element in $Y_2\backslash Y_1$ such that $\rho(z_1,Y_0)=2$, $\rho(z_1,Y_1)=1$; $w_1$ is an element in $\overline Y_2$ with $\|z_1-w_1\|=1$.
    \item[(ii)] $f_1:~Y_2\rightarrow \mathbb R$ is a non-zero linear functional such that 
    $$
    f_1(z_1)=\lim\limits_{a\to\infty}\left(a-\frac{\rho(z_1-aw_1,Y_0))}{\rho(w_1,Y_0)}\right).
    $$
    \item[(iii)] $\mu_n^{(1)}$ is some value in $[v_n^{(1)},u_n^{(1)}]=\left[1,1+\frac{\tau_n}{2d_1}\right]$.
    \item[(iv)] We have 
$$
\rho(q_{1,n},Y_{0})=c^{(1)}\left(1+\frac{\tau_n}{2d_1}\right),~\rho(q_{1,n},Y_1)=1,
$$
where $c^{(1)}=1/|1-2f_1(z_1)|$ 
and there exists another constant $c_1>0$ such that
$$
\|q_{1,m}-q_{1,n}\|\le c_1\frac{\tau_m}{2d_1},~\mbox{for all $n\ge m\ge 1$}.
$$
\end{enumerate}
\item For $j=2$, using the condition (\ref{Yn}), we can find $y_2\in Y_3\backslash Y_2$ such that
$$
\rho\left(\frac{y_2}{\|y_2\|},Y_2\right)=1.
$$
Let 
$$
z_2'=\frac{y_2}{\|y_2\|}+z_1,
$$
then $z_2'\in Y_3\backslash Y_2$ and
\begin{equation}
\label{upper:z2}
\|z_2'-z_1\|=\rho(z_2'-z_1,Y_2)=\rho(z_2',Y_2)=1.
\end{equation}
Next we construct $z_2$ as well as $q_{2,n}$. To this end, let
$$
\alpha_{min}=-1/2;~\alpha_{max}=1/2.
$$
By using the triangle inequality and (\ref{upper:z2}) we have
\begin{equation}
\label{alpha_min}
\|z_2'+\alpha_{min}z_1\|\le \|z_2'-z_1\|+|1+\alpha_{min}|\|z_1\|=1+2|1+\alpha_{min}|=2
\end{equation}
and
\begin{equation}
    \label{lower:z2}
    \|z_2'+\alpha_{max} z_1\|\ge|1+\alpha_{max}|\|z_1\|-\|z_2'-z_1\|=2|1+\alpha_{max}|-1=2.
    \end{equation}
Then the continuity of the mapping $\lambda\mapsto \rho(z_2'+\lambda z_1,Y_1)$,  (\ref{alpha_min}), (\ref{lower:z2}) and the intermediate value theorem imply the existence of a real number $\alpha_2\in[-1/2,1/2]$ such that
$$
\rho(z_2'+\alpha_2z_1,Y_0)=2.
$$
Let
$$
z_2=z_2'+\alpha_2z_1,~z_1'=(\alpha_2+1)z_1,
$$ then  $z_2\in Y_3\backslash Y_2$ satisfies $\rho(z_2,Y_0)=2$, $\rho(z_2,Y_2)=1$; and $z_1'$ satisfies $z_1'\in Y_2\backslash Y_0$ and $\|z_2-z_1'\|=1$. We thus can apply Lemma \ref{lemma2} (by taking $(z,w)=(z_2,z_1')$ in the lemma) to obtain, the existence of a sequence $\{q_{2,n}\}_{n\ge 2}\subset Y_{3}\backslash Y_2$ given by:
$$
q_{2,n}=(z_2-f_2(z_2)z_1')+\mu_n^{(2)}f_2(z_2)z_1',
$$
where
\renewcommand{\labelenumi}{\alph{enumi})}
\begin{enumerate}
    \item[(i)] The non-zero linear functional
     $f_2:~Y_3\rightarrow \mathbb R$ is such that 
     $$
     f_2(z_2)=\lim\limits_{a\to\infty}\left(a-\frac{\rho(z_2-az_1',Y_0)}{\rho(z_1',Y_0)}\right)\ge 0.
     $$
    \item[(ii)] $\mu_n^{(2)}$ is some value in $[v_n^{(2)},u_n^{(2)}]=\left[1,1+\frac{\tau_n}{2^2d_2}\right]$ such that
    $$
\rho(q_{2,n},Y_0)=c^{(2)}\left(1+\frac{\tau_n}{2^2d_2}\right),~\rho(q_{2,n},Y_2)=1
$$
for some constant $c^{(2)}>0$.
\item[(iii)] There is a constant $c_2>0$ such that
$$
\|q_{2,m}-q_{2,n}\|\le c_2\frac{\tau_m}{2^2d_2},~\mbox{for all $n\ge m\ge 2$}.
$$
Therefore, by using Lemma \ref{lemma1'},  we get the existence of a non-zero linear functional $f_{1,n}:~Y_{3}\longmapsto\mathbb R$ such that
$$
Y_1\subset \ker f_{1,n},~\|f_{1,n}\|=\frac{1}{\rho(q_{1,n},Y_1)}=f_{1,n}(q_{1,n})=1
$$
and
\begin{eqnarray*}
f_{1,n}(q_{2,n})&\ge&(\mu_n^{(2)}-1)f_2(z_2)-\frac{\rho\left(q_{2,n}-((\mu_n^{(2)}-1)f_2(z_2))q_{1,n},Y_1\right)}{\rho(q_{1,n},Y_1)}.
\end{eqnarray*}
Since
\begin{eqnarray*}
&&\rho\left(q_{2,n}-(\mu_n^{(2)}-1)f_2(z_2)q_{1,n},Y_1\right)
\\
&&\le \rho\left(q_{2,n},Y_1\right)+(\mu_n^{(2)}-1)f_2(z_2)\rho\left(q_{1,n},Y_1\right)\\
&&=\rho\left(q_{2,n},Y_1\right)+(\mu_n^{(2)}-1)f_2(z_2).
\end{eqnarray*}
\end{enumerate}
\end{itemize}
Continuing with  the steps above we can obtain that:
for $j=1,\ldots,n-1$, there exists a sequence $\{q_{j,n}\}_{n\ge j}\subset Y_{j+1}\backslash Y_j$ such that
\begin{equation}
\label{qjnY}
\rho(q_{j,n},Y_{0})=1+\frac{\tau_n}{2^jd_j},~\rho(q_{j,n},Y_j)=1;
\end{equation}
\begin{equation}
\label{diffq}
\|q_{j,m}-q_{j,n}\|\le 4\frac{\tau_m}{2^jd_j},~\mbox{for all $n\ge m\ge j$}.
\end{equation}
Also there is a non-zero linear functional $f_{j,n}:~Y_{j+2}\longmapsto\mathbb R$ such that
$$
Y_j\subset \ker f_{j,n},~\|f_{j,n}\|=f_{j,n}(q_{j,n})=1
$$
and
\begin{equation}
\label{f}
f_{j,n}(q_{j+1,n})\in\left[-1,\frac{\tau_n}{2^{j}d_{j+1}}-1\right].
\end{equation}
Now we fix $n\ge1$. Take $\lambda_{n,n}:=d_n$. We see clearly from (\ref{qjnY}) that
\begin{equation}
\label{identqn}
\rho(\lambda_{n,n}q_{n,n},Y_n)=\lambda_{n,n}\rho(q_{n,n},Y_n)=d_n.
\end{equation}
 Then, first by using (\ref{qjnY}) and (\ref{f}), we obtain
\begin{eqnarray}
\label{upperqn}
&&\rho(\lambda_{n,n}q_{n,n},Y_{n-1})=\lambda_{n,n}\rho(q_{n,n},Y_{n-1})\le \lambda_{n,n}\rho(q_{n,n},Y_0)\nonumber\\
&&=d_n\left(1+\frac{\tau_n}{2^nd_n}\right)\le  d_n+\tau_n\le d_{n-1};
\end{eqnarray}
and, by the properties of $f_{n-1,n}$ in (\ref{f}), we see that
\begin{eqnarray}
\label{lowerqn}
&&\rho(\lambda_{n,n}q_{n,n}-(d_{n-1}+d_nf_{n-1,n}(q_{n,n}))q_{n-1,n},Y_{n-1})\nonumber\\
&&\ge |f_{n-1,n}(\lambda_{n,n}q_{n,n}-(d_{n-1}+d_nf_{n-1,n}(q_{n,n}))q_{n-1,n})|\nonumber\\
&&= -d_nf_{n-1,n}(q_{n,n})+(d_{n-1}+d_nf_{n-1,n}(q_{n,n}))f_{n-1,n}(q_{n-1,n})\nonumber\\
&&=d_{n-1}+d_n\left(f_{n-1,n}(q_{n,n})-f_{n-1,n}(q_{n,n})\right)= d_{n-1}.
\end{eqnarray}
The mapping
$$
\lambda\longmapsto\rho(\lambda_{n,n}q_{n,n}+\lambda q_{n-1,n},Y_{n-1})
$$
 is continuous, through which the image of $[-(d_{n-1}+d_nf_{n-1,n}(q_{n,n})),0]$ contains $d_{n-1}$. Then,  by (\ref{upperqn}), (\ref{lowerqn}) and the intermediate value theorem, there exists $\lambda_{n-1,n}\in [-(d_{n-1}+d_nf_{n-1,n}(q_{n,n})),0]$ such that
\begin{equation}
\label{rhodn}
\rho(\lambda_{n,n}q_{n,n}+\lambda_{n-1,n}q_{n-1,n},Y_{n-1})= d_{n-1}.
\end{equation}
Furthermore, by the fact that $q_{n-1,n}\in Y_n$ and (\ref{identqn}), we have
$$
\rho(\lambda_{n,n}q_{n,n}+\lambda_{n-1,n}q_{n-1,n},Y_{n})= \rho(\lambda_{n,n}q_{n,n},Y_{n})= d_{n}.
$$
For any $2\le k\le n-1$, assume that we can find the real numbers
\begin{eqnarray}
\label{int-lambda}
&&\lambda_{n-1,n}\in[-(d_{n-1}+d_{n}f_{n-1,n}(q_{n,n})),0],\nonumber\\
&&~\ldots,\ldots,\ldots~\nonumber\\
&&\lambda_{k,n}\in[-(d_{k}+d_{k+1}f_{k,n}(q_{k+1,n})),0]
\end{eqnarray}
such that
$$
\rho(\lambda_{n,n}q_{n,n}+\lambda_{n-1,n}q_{n-1,n}+\ldots+\lambda_{k,n}q_{k,n},Y_m)=d_m,~\mbox{for}~m=k,\ldots,n.
$$
Let $
z_{k,n}=\lambda_{n,n}q_{n,n}+\ldots+\lambda_{k,n}q_{k,n}$.
Then, on one hand, by using the triangle inequality, (\ref{int-lambda}), (\ref{qjnY}), (\ref{f}) and  the fact that $k\ge2$, we obtain
\begin{eqnarray}
\label{zY}
&&\rho(z_{k,n},Y_{k-1})\le\sum_{j=k}^n|\lambda_{j,n}|\rho(q_{j,n},Y_{k-1})\nonumber\\
&&\le|\lambda_{n,n}|\rho(q_{n,n},Y_{0})+\sum_{j=k}^{n-1}|\lambda_{j,n}|\rho(q_{j,n},Y_{0})\nonumber\\
&&\le d_n\left(1+\frac{\tau_{n}}{2^nd_n}\right)+\sum_{j=k}^{n-1}\left(d_{j}+d_{j+1}f_{j,n}(q_{j+1,n})\right)\left(1+\frac{\tau_{n}}{2^jd_j}\right)\nonumber\\
&&\le d_n\left(1+\frac{\tau_{n}}{2^nd_n}\right)+\sum_{j=k}^{n-1}\left(d_{j}+d_{j+1}\left(\frac{\tau_n}{2^{j}d_{j+1}}-1\right)\right)\left(1+\frac{\tau_{n}}{2^jd_j}\right)\nonumber\\
&&\le \left(d_n+\sum_{j=k}^{n-1}(d_{j}-d_{j+1})\right)+\sum_{j=2}^{n}\frac{\tau_n(2d_j+d_{j+1}(\tau_n/(2^{j}d_{j+1})-1)}{2^jd_j}\nonumber\\
&&\le d_k+\sum_{j=2}^{n}\frac{\tau_n(2d_j)}{2^{j}d_j}= d_k+\tau_n\sum_{j=2}^n\frac{1}{2^{j-1}}\nonumber\\
&&\le d_k+\tau_n\sum_{j=2}^\infty\frac{1}{2^{j-1}}= d_k+\tau_n\le d_{k-1}.
\end{eqnarray}
On the other hand, by using the properties of $f_{k-1,n}$ in (\ref{f}), we obtain
\begin{equation}
\label{bound:z}
\rho(z_{k,n}-(d_{k-1}+d_kf_{k-1,n}(q_{k-1,n})q_{k-1,n},Y_{k-1})\ge d_{k-1}f_{k-1,n}(q_{k-1,n})= d_{k-1}.
\end{equation}
Therefore, in view of (\ref{zY}), (\ref{bound:z}) and the intermediate value theorem, there is $\lambda_{k-1,n}\in [-(d_{k-1}+f_{k-1,n}(q_{k-1,n})),0]$ such that
$$
\rho(z_{k,n}+\lambda_{k-1,n}q_{k-1,n},Y_{k-1})=d_{k-1}.
$$
Furthermore,
$$
\rho(z_{k,n}+\lambda_{k-1,n}q_{k-1,n},Y_{m})=\rho(z_{k,n},Y_m)=d_{m}~\mbox{for}~m=k,\ldots,n.
$$
Continuing this procedure until $k=1$ is included, we obtain the element
$$
x_{n,n}=\lambda_{n,n}q_{n,n}+\ldots+\lambda_{1,n}q_{1,n},
$$
for which $\rho(x_{n,n},Y_k)=d_k$ and
\begin{equation}
\label{bound_lambdakn}
|\lambda_{k,n}|\left\{
\begin{array}{ll}
\le d_k-d_{k+1}(1-2^{-k})& \mbox{for all $n\ge2$, $k\le n-1$,}\\
=d_n&\mbox{for $n\ge1$, $k=n$.}
\end{array}\right.
 \end{equation}
We see from (\ref{bound_lambdakn}) that for each $k\ge1$, the sequence $\{\lambda_{k,n}\}_{n\ge k}$ is bounded by $d_k$. Then using the usual diagonalization process, we choose a sequence $\Lambda$ of indices $n$ such that, for all $k\ge1$, $\lambda_{k,n}$ converges to the limit $\lambda_k$ as $n\rightarrow\infty$, $n\in\Lambda$. We then claim that $|\lambda_k|\le d_k$ and the limit $\lim\limits_{n\rightarrow\infty}\sum\limits_{k=1}^n\lambda_kq_{k,n}$ exists in $X$. Indeed, by using the triangle inequality, we have for $m\le n$,
\begin{eqnarray}
\label{dfflambdaq}
&&\Big\|\sum_{k=1}^n\lambda_{k}q_{k,n}-\sum_{k=1}^m\lambda_{k}q_{k,m}\Big\|=\Big\|\sum_{k=1}^m\lambda_{k}(q_{k,n}-q_{k,m})+\sum_{k={m+1}}^n\lambda_{k}q_{k,n}\Big\|\nonumber\\
&&\le \sum_{k=1}^m|\lambda_{k}|\|q_{k,n}-q_{k,m}\|+\sum_{k={m+1}}^n|\lambda_{k}|\|q_{k,n}\|.
\end{eqnarray}
First using $|\lambda_k|\le d_k$ and (\ref{diffq}), we obtain
\begin{equation}
\label{dfflambdaq1}
\sum_{k=1}^m|\lambda_{k}|\|q_{k,n}-q_{k,m}\|\le 4\tau_m\sum_{k=1}^m\frac{d_k}{2^kd_k}\xrightarrow[n,m\rightarrow\infty]{}0,
\end{equation}
and  from (\ref{zY}) we see
\begin{equation}
\label{dfflambdaq2}
\sum_{k=m}^n|\lambda_k|\|q_{k,n}\|\le d_{m-1}\xrightarrow[n,m\rightarrow\infty]{}0.
\end{equation}
Combining (\ref{dfflambdaq}), (\ref{dfflambdaq1}) and (\ref{dfflambdaq2}) yields
$$
\left\|\sum_{k=1}^n\lambda_{k}q_{k,n}-\sum_{k=1}^m\lambda_{k}q_{k,m}\right\|\xrightarrow[n,m\rightarrow\infty]{}0,
$$
i.e., $\Big\{\sum\limits_{k=1}^n\lambda_kq_{k,n}\Big\}_{n\ge1}$ is a Cauchy sequence in the Banach space $X$, therefore it has a limit in $X$.
Furthermore, we claim that the element
$$
x:=\lim_{n\rightarrow\infty}\sum_{k=1}^n\lambda_kq_{k,n}
$$
is the limit of the sequence $\{x_{n,n}\}_{n\in\Lambda}$ as $n\to\infty$. By using the fact that
$\|q_{k,n}\|\le 2$, (\ref{bound_lambdakn}) and (\ref{dfflambdaq2}), we obtain
\begin{eqnarray*}
&&\|x-x_{n,n}\|\le \|x-x_n\|+\|x_n-x_{n,n}\|\le \|x-x_n\|+\sum_{k=1}^n|\lambda_{k,n}-\lambda_k|\|q_{k,n}\|\\
&&\le \|x-x_n\|+\sum_{k=1}^{N}|\lambda_{k,n}-\lambda_k|\|q_{k,n}\|+\sum_{k=N+1}^n(|\lambda_{k,n}|+|\lambda_k|)\|q_{k,n}\|\\
&&\le \|x-x_n\|+2\sum_{k=1}^{N}|\lambda_{k,n}-\lambda_k|+d_{n}\|q_{n,n}\|\\
&&+2\sum_{k=N+1}^{n-1}\left(d_k-d_{k+1}+d_{k+1}2^{-k}\right)+\sum_{k=N+1}^n|\lambda_k|\|q_{k,n}\|\\
&&\le \|x-x_n\|+2 N\max_{1\le k\le N,N<n}|\lambda_{k,n}-\lambda_k|+2d_n\\
&&+2\left(d_{N+1}-d_n+d_{N+1}\sum_{k=N+1}^{\infty}2^{-k}\right)+\sum_{k=N+1}^n|\lambda_k|\|q_{k,n}\|\xrightarrow[n\rightarrow\infty,n\in\Lambda,N\rightarrow\infty]{}0.
\end{eqnarray*}
Finally, by the continuity of $x\longmapsto\rho(x,Y_k)$, we obtain
$$
\rho(x,Y_k)=\lim_{n\rightarrow\infty,n\in\Lambda}\rho(x_{n,n},Y_k)=d_k,~\mbox{for all}~k\ge1.
$$
Now we prove the theorem in the following two other particular cases:
\begin{enumerate}
\item If $Y_1=\{0\}$, the problem can be easily converted to the case $Y_1\neq\{0\}$:
\begin{itemize}
\item If $d_1>d_2$, then having constructed an element $z$ with $\rho(z,Y_n)=d_n$ for all $n\ge 2$, we can use Lemma \ref{lemma1} to construct an element $x$ with $\rho(x,Y_k)=d_k$ for $k=1,2$ and such that $x-\lambda z\in Y_{2}$ for some $\lambda>0$. But then observe that
$$
d_{2}=\rho(x,Y_{2})=\rho((x-\lambda z)+\lambda z,Y_2)=\rho(\lambda z,Y_{2})=\lambda d_{2};
$$
therefore, $\lambda=1$ and
$$\rho(x,Y_n)=\rho(z,Y_n)=d_n~\mbox{ for all $n\ge 2$}.
$$
This together with the fact that $\rho(x,Y_1)=d_1$ entails $\rho(x,Y_n)=d_n$ for all $n\ge1$.
\item If $d_1=d_2$, then it suffices to take
$$
q_{1,n}=z_1
$$
with $z_1\in Y_2\backslash Y_1$ satisfying $\|z_1\|=1$. The constructions of
$$q_{2,n},~q_{3,n},\ldots$$ are similar to the above proof.
\end{itemize}
\item If $d_n=0$ starting from some $n$, then the desired element exists by applying the above proof to the setting $X=Y_n$  where the subspaces $Y_1\subset Y_2\subset\ldots\subset Y_{n-1}$  are in $X$.
\end{enumerate}
\end{proof}

\end{document}
